% Technical Log Template
%
% The Technical Log is a comprehensive record of everything the agent did. It should be
% exhaustive—this is the authoritative documentation that enables complete understanding
% and reproduction of the analysis.
%
% LaTeX Template
\documentclass[12pt,modern,linenumbers]{aastex7}

\usepackage{listings}
\usepackage{tcolorbox}
\usepackage{longtable}
\usepackage{xcolor}

% Code listing style
\lstset{
    language=Python,
    basicstyle=\ttfamily\small,
    keywordstyle=\color{blue},
    commentstyle=\color{gray},
    stringstyle=\color{red!70!black},
    frame=single,
    breaklines=true,
    numbers=left,
    numberstyle=\tiny\color{gray},
    xleftmargin=2em
}

% Output box
\newtcolorbox{codeoutput}{
    colback=black!5!white,
    colframe=black!50!white,
    title=Output,
    fonttitle=\bfseries\small
}

% Decision box
\newtcolorbox{decision}[1]{
    colback=blue!5!white,
    colframe=blue!75!black,
    fonttitle=\bfseries,
    title={Decision: #1}
}

% Warning box
\newtcolorbox{warning}{
    colback=red!5!white,
    colframe=red!75!black,
    title=Warning,
    fonttitle=\bfseries
}

\begin{document}

\title{Technical Log: [Analysis Title]}

\author{[Agent Name / Analysis System]}
\affiliation{Generated automatically from execution trace}
\email{N/A}

\begin{abstract}
This technical log documents the complete execution trace of [brief description 
of analysis]. It includes all data operations, model training steps, decision 
points, and their outcomes. The log is intended to enable complete reproduction 
and evaluation of the analysis methodology.

\textbf{Execution Period}: [Start date] -- [End date]\\
\textbf{Total Runtime}: [X hours]\\
\textbf{Final Status}: [Completed / Partial / Failed]
\end{abstract}

%% ============================================================================
\section{Execution Summary}
%% ============================================================================

\subsection{Objective}
[What was the agent trying to accomplish?]

\subsection{Outcome}
[What was achieved? What metrics were obtained?]

\subsection{Key Findings}
\begin{enumerate}
    \item [Finding 1]
    \item [Finding 2]
    \item [Finding 3]
\end{enumerate}

\subsection{Known Issues}
\begin{itemize}
    \item [Issue 1: description and impact]
    \item [Issue 2: description and impact]
\end{itemize}

%% ============================================================================
\section{Environment}
%% ============================================================================

\subsection{Hardware}
\begin{itemize}
    \item CPU: [Model, cores]
    \item GPU: [Model, memory]
    \item RAM: [Amount]
    \item Storage: [Type, capacity]
\end{itemize}

\subsection{Software}
\begin{lstlisting}
# Key dependencies (from environment.yml)
python=3.10.12
numpy=1.24.3
pytorch=2.0.1+cu118
astropy=5.3.1
scikit-learn=1.3.0
\end{lstlisting}

\subsection{Reproducibility Information}
\begin{itemize}
    \item Random seed: 42
    \item CUDA deterministic: True
    \item PyTorch version: 2.0.1
\end{itemize}

%% ============================================================================
\section{Data Provenance}
%% ============================================================================

\subsection{Source Data}

\begin{deluxetable*}{llllc}
\tablecaption{Data Sources\label{tab:data_sources}}
\tablehead{
\colhead{Dataset} & \colhead{Source} & \colhead{Access Date} & 
\colhead{Format} & \colhead{Size}
}
\startdata
DESI DR1 Catalog & legacysurvey.org/dr9 & 2024-01-15 & FITS & 12.3 GB \\
Galaxy Images & cutout service & 2024-01-16 & FITS & 45.2 GB \\
\enddata
\end{deluxetable*}

\subsection{Selection Criteria}

The following query was used to select the sample:

\begin{lstlisting}[language=SQL]
SELECT objid, ra, dec, flux_g, flux_r, flux_z, z_spec
FROM dr9_south
WHERE r_mag < 21.2
  AND stellar_mass < POWER(10, 9.25)
  AND z_spec IS NOT NULL
  AND z_spec BETWEEN -0.01 AND 1.5
-- Result: 249,999 galaxies
\end{lstlisting}

\subsection{Data Quality Assessment}

\begin{lstlisting}
# Quality check results
nan_fraction = 0.0%
inf_fraction = 0.0%
duplicate_fraction = 0.0%
missing_images = 0
\end{lstlisting}

%% ============================================================================
\section{Chronological Execution Trace}
%% ============================================================================

% Repeat this pattern for each major step

\subsection{Step 1: Data Loading}
\textbf{Timestamp}: 2024-01-15 10:23:45 UTC\\
\textbf{Duration}: 3 minutes 22 seconds

\begin{lstlisting}
# Load DESI catalog
from astropy.io import fits

catalog = fits.getdata('data/raw/desi_dr1.fits')
print(f"Loaded {len(catalog)} galaxies")
\end{lstlisting}

\begin{codeoutput}
Loaded 1,234,567 galaxies
\end{codeoutput}

\begin{decision}{D1: Magnitude Cut}
\textbf{Choice}: $r < 21.2$

\textbf{Alternatives considered}:
\begin{itemize}
    \item $r < 20.5$: Higher quality but 40\% smaller sample
    \item $r < 22.0$: Larger sample but significantly more noise
    \item $r < 21.5$: Compromise option
\end{itemize}

\textbf{Rationale}: The $r < 21.2$ cut provides adequate signal-to-noise for 
morphological measurements while maintaining a statistically significant sample 
of low-redshift dwarf galaxies.

\textbf{Sensitivity test}:
\begin{lstlisting}
# Test different magnitude cuts
for mag_cut in [20.5, 21.0, 21.2, 21.5, 22.0]:
    n_galaxies = np.sum(catalog['r_mag'] < mag_cut)
    n_lowz = np.sum((catalog['r_mag'] < mag_cut) & 
                    (catalog['z_spec'] < 0.03))
    print(f"r < {mag_cut}: {n_galaxies} total, {n_lowz} low-z")
\end{lstlisting}

\begin{codeoutput}
r < 20.5: 148,234 total, 8,921 low-z
r < 21.0: 198,456 total, 12,345 low-z
r < 21.2: 249,999 total, 15,102 low-z
r < 21.5: 312,678 total, 17,234 low-z
r < 22.0: 456,789 total, 19,567 low-z
\end{codeoutput}
\end{decision}

\subsection{Step 2: Data Preprocessing}
\textbf{Timestamp}: 2024-01-15 10:27:07 UTC\\
\textbf{Duration}: 15 minutes 48 seconds

[Continue with similar structure for each step...]

%% ============================================================================
\section{Model Training Log}
%% ============================================================================

\subsection{Architecture}

\begin{lstlisting}
MultiTaskCNN(
  backbone: ResNet-18 (11.2M parameters)
  regression_head: Linear(512 -> 1)
  binary_head: Linear(512 -> 1) + Sigmoid
  multiclass_head: Linear(512 -> 3) + Softmax
)
Total parameters: 11,178,945
\end{lstlisting}

\subsection{Training Configuration}

\begin{deluxetable}{ll}
\tablecaption{Training Hyperparameters\label{tab:hyperparams}}
\tablehead{
\colhead{Parameter} & \colhead{Value}
}
\startdata
Optimizer & Adam \\
Learning rate & $10^{-3}$ \\
Batch size & 256 \\
Max epochs & 50 \\
Early stopping patience & 10 \\
Loss weights (λ$_1$, λ$_2$) & (0.5, 0.5) \\
\enddata
\end{deluxetable}

\subsection{Training Progress}

\begin{lstlisting}
Epoch  1/50: train_loss=0.0452, val_loss=0.0521, val_nmad=0.0412
Epoch  2/50: train_loss=0.0389, val_loss=0.0445, val_nmad=0.0356
...
Epoch 23/50: train_loss=0.0198, val_loss=0.0267, val_nmad=0.0265
Epoch 24/50: train_loss=0.0195, val_loss=0.0270, val_nmad=0.0268
Early stopping triggered at epoch 24
Best model saved at epoch 23 with val_nmad=0.0265
\end{lstlisting}

%% ============================================================================
\section{Decision Registry}
%% ============================================================================

\begin{deluxetable*}{clllc}
\tablecaption{Complete Decision Registry\label{tab:decision_registry}}
\tablehead{
\colhead{ID} & \colhead{Decision} & \colhead{Alternatives} & 
\colhead{Rationale} & \colhead{Sensitivity}
}
\startdata
D1 & $r < 21.2$ magnitude cut & 20.5, 21.0, 21.5, 22.0 & S/N vs sample size & 40\% sample change \\
D2 & Per-image normalization & Global, per-band & Varying conditions & 15\% σ$_{\rm NMAD}$ \\
D3 & ResNet-18 backbone & VGG, EfficientNet & Speed/accuracy & ±2\% σ$_{\rm NMAD}$ \\
D4 & Huber loss (δ=1.0) & MSE, MAE & Outlier robustness & 8\% outlier rate \\
D5 & Multi-task learning & Single-task & Low-z performance & 57\% low-z outliers \\
D6 & 70/15/15 split & 80/10/10 & Validation size & <1\% effect \\
D7 & Batch size 256 & 64, 128, 512 & GPU memory & <1\% effect \\
D8 & Learning rate $10^{-3}$ & $10^{-2}$, $10^{-4}$ & Convergence speed & 2× training time \\
D9 & λ = 0.5 auxiliary weights & 0.1, 0.3, 1.0 & Task balance & 5\% low-z outliers \\
D10 & 128×128 pixel images & 64×64, 256×256 & Resolution vs compute & 26\% σ$_{\rm NMAD}$ \\
\enddata
\end{deluxetable*}

%% ============================================================================
\section{Failure Log}
%% ============================================================================

\subsection{Failed Approaches}

\subsubsection{Attempt 1: Feature Engineering with XGBoost}
\textbf{Date}: 2024-01-16\\
\textbf{Description}: Trained XGBoost on 64 handcrafted morphological features.

\begin{warning}
This approach failed with severe overfitting: validation σ$_{\rm NMAD}$ = 0.03 
but test σ$_{\rm NMAD}$ = 0.12 (4× worse).
\end{warning}

\textbf{Diagnosis}: Handcrafted features were too specific to training set 
characteristics and did not generalize.

\textbf{Resolution}: Switched to end-to-end CNN learning.

\subsubsection{Attempt 2: Scale-Normalized Images}
\textbf{Date}: 2024-01-17\\
\textbf{Description}: Cropped images to 3×R50 to remove angular size information.

\begin{warning}
Performance degraded by 4.1\% (MAE: 0.0336 → 0.0350). Scale and morphology are 
physically coupled for resolved galaxies.
\end{warning}

\textbf{Diagnosis}: Angular size is a valid distance indicator for resolved 
sources; removing it discards useful information.

\textbf{Resolution}: Retained full 128×128 images with scale information.

%% ============================================================================
\section{Reproducibility Manifest}
%% ============================================================================

\subsection{Files Required}

\begin{deluxetable}{llc}
\tablecaption{Required Files for Reproduction\label{tab:manifest}}
\tablehead{
\colhead{File} & \colhead{Description} & \colhead{Size}
}
\startdata
src/scripts/01\_download.py & Data download script & 2 KB \\
src/scripts/02\_preprocess.py & Preprocessing pipeline & 8 KB \\
src/scripts/03\_train.py & Model training & 15 KB \\
src/scripts/04\_evaluate.py & Evaluation script & 6 KB \\
environment.yml & Conda environment & 1 KB \\
Snakefile & Build automation & 3 KB \\
models/multitask\_cnn\_best.pth & Trained model & 134 MB \\
\enddata
\end{deluxetable}

\subsection{Reproduction Commands}

\begin{lstlisting}[language=bash]
# From a fresh clone:
git clone https://github.com/[repo].git
cd [repo]
conda env create -f environment.yml
conda activate photoz

# Full reproduction (including training):
snakemake --cores 4

# Using pre-trained model (faster):
snakemake --cores 4 --config use_pretrained=True
\end{lstlisting}

\subsection{Expected Outputs}

\begin{itemize}
    \item \texttt{outputs/predictions.csv}: Test set predictions (37,500 rows)
    \item \texttt{outputs/metrics.json}: Performance metrics
    \item \texttt{figures/*.pdf}: All figures
    \item \texttt{outputs/report.pdf}: Summary report
\end{itemize}

%% ============================================================================
\section{Appendix: Raw Logs}
%% ============================================================================

[Include verbatim agent output logs if available]

\end{document}
%
% Key Sections Explained
%
% Execution Summary
% - Quick overview for readers who need the TL;DR
% - What was attempted, what succeeded, what failed
% - Key numbers and findings
%
% Environment
% - Complete specification for reproduction
% - Includes hardware (affects training time)
% - Software versions matter for reproducibility
%
% Data Provenance
% - Every data source with URL and access date
% - Exact queries used for selection
% - Quality checks performed and results
%
% Chronological Trace
% - Step-by-step record of what was done
% - Includes actual code and outputs
% - Decision boxes for each choice point
%
% Decision Registry
% - Complete table of all decisions
% - Enables sensitivity analysis
% - Quick reference for methodology choices
%
% Failure Log
% - Documents what didn't work
% - Explains why approaches failed
% - Shows diagnostic reasoning
% - This is often the most valuable section for future researchers!
%
% Reproducibility Manifest
% - Complete list of files needed
% - Clear reproduction commands
% - Expected outputs for verification
%
% Tips for Good Technical Logs
%
% 1. **Be exhaustive** - Include everything, even if it seems obvious
% 2. **Timestamp everything** - Enables reconstruction of sequence
% 3. **Show actual outputs** - Not paraphrased descriptions
% 4. **Document failures** - They're informative
% 5. **Include sensitivity tests** - Show robustness of choices
% 6. **Use consistent formatting** - Decision boxes, code blocks, etc.
% 7. **Link to code** - Reference specific scripts and line numbers
% 8. **Make it searchable** - Use clear section headings and labels